\chapter{Multi-agent reinforcement learning}
\label{ch:multiagent}

Pacman shares its maze with ghosts. In this course they follow fixed rules,
so from Pacman's point of view they are part of the environment, and the
problem is an MDP. If the ghosts learned too, each agent would face an
environment that changes as the others learn: the problem would become a
\emph{game}. Multi-agent RL covers competitive games (chess, Go, poker,
StarCraft), cooperative teams (robots, traffic lights, network routing)
and mixed situations (markets, negotiation).

\section{Markov games}

\begin{definition}{Markov game}{markov-game}
  A Markov (or stochastic) game with $n$ agents is a tuple $(\States,
  \Actions_1, \dots, \Actions_n, p, r_1, \dots, r_n, \gamma)$: in state $s$,
  every agent $i$ chooses an action $a_i$, the joint action $\mathbf{a} =
  (a_1, \dots, a_n)$ determines the transition $p(s' \mid s, \mathbf{a})$,
  and agent $i$ receives $r_i(s, \mathbf{a})$. Each agent maximises its own
  expected discounted return.
\end{definition}

With one agent it is an MDP; with one state it is a repeated matrix game.
The game is \emph{fully cooperative} if all the rewards are equal,
\emph{zero-sum} if $n = 2$ and $r_1 = -r_2$. In general there is no single
optimal policy: the best policy of an agent depends on those of the others.
A joint policy $(\pi_1, \dots, \pi_n)$ is a \emph{Nash equilibrium} if no
agent can increase its return by changing its policy alone; every finite
discounted Markov game has a Nash equilibrium in stationary stochastic
policies (Fink, 1964).

\section{Two-player zero-sum games}

Zero-sum games are the multi-agent case where the single-agent theory
extends most completely. In a matrix game with payoff matrix $A$ (the
first player receives $x^\top A y$ for mixed strategies $x, y$), von
Neumann's minimax theorem states that $\max_x \min_y x^\top A y = \min_y
\max_x x^\top A y$; this common \emph{value} $\mathrm{val}(A)$ is computed
by linear programming. Shapley (1953) extended the Bellman equations: the
value of a zero-sum Markov game satisfies
\begin{equation}
  v(s) = \mathrm{val}\Bigl[ r(s, a_1, a_2) + \gamma \sum_{s'} p(s' \mid s,
  a_1, a_2)\, v(s') \Bigr]_{a_1, a_2},
\end{equation}
the value of the matrix game indexed by the two actions.

\begin{theorem}{Shapley}{shapley}
  The operator $\T$ defined by the right-hand side is a $\gamma$-contraction
  in the sup norm. Its unique fixed point is the value of the game, and the
  optimal mixed strategies of the matrix games at the fixed point form a
  stationary equilibrium.
\end{theorem}

\begin{proof}
  We prove the contraction. First, $\mathrm{val}$ is non-expansive: if
  $\abs{A_{jk} - B_{jk}} \leq \delta$ for all $j, k$, then $x^\top A y \leq
  x^\top B y + \delta$ for all mixed strategies, since $x^\top (A - B) y$
  is an average of entries of $A - B$; taking $\min_y$ then $\max_x$ gives
  $\mathrm{val}(A) \leq \mathrm{val}(B) + \delta$, and symmetrically. For two
  functions $v, w$, the matrices defining $(\T v)(s)$ and $(\T w)(s)$
  differ entrywise by $\gamma \sum_{s'} p(s' \mid s, a_1, a_2) (v(s') -
  w(s'))$, at most $\gamma \norm{v - w}_\infty$ in absolute value. Hence
  $\norm{\T v - \T w}_\infty \leq \gamma \norm{v - w}_\infty$. The rest
  follows as for MDPs (\cref{thm:bellman-opt}).
\end{proof}

Minimax-Q (Littman, 1994) is the corresponding Q-learning: it learns
$Q(s, a_1, a_2)$ with the target $r + \gamma\, \mathrm{val}[Q(s', \cdot,
\cdot)]$, solving a small linear program at each update, and converges to
the value of the game under the same conditions as Q-learning. Nash-Q
extends the idea to general-sum games, but needs restrictive assumptions:
equilibria are no longer unique or interchangeable.

\paragraph{Self-play.} In symmetric games, an agent can learn by playing
against copies of itself. TD-Gammon (Tesauro, 1995), AlphaZero
(\cref{ch:model-based}) and OpenAI Five learnt this way. Pure self-play
may cycle (as in rock--paper--scissors, where each best response is beaten
by the next) or overfit to its own style. \emph{Fictitious play} best
responds to the \emph{average} of the past opponents, and converges to
equilibrium in zero-sum games; population methods keep a league of past
and exploiting agents (AlphaStar, Vinyals et al., 2019) or compute
equilibria of the meta-game between policies (PSRO, Lanctot et al.,
2017). Imperfect-information games such as poker require different tools
(counterfactual regret minimisation, which beat professional players at
heads-up and six-player no-limit hold'em).

\section{Learning in cooperative teams}

\paragraph{Independent learners.} The simplest approach lets each agent run
its own single-agent algorithm (independent Q-learning, or independent PPO),
treating the others as part of the environment. It often works, but the
environment of each agent is non-stationary while the others learn, so the
convergence guarantees are lost, and the replay memory of DQN holds
transitions generated by obsolete policies of the others.

\paragraph{Centralised training, decentralised execution.} In simulation,
training can use information that will not be available at execution: the
global state and the actions of all the agents. MADDPG (Lowe et al., 2017)
gives each agent a centralised critic $Q_i(s, a_1, \dots, a_n)$ and a
decentralised actor $\mu_i(o_i)$ of its own observation, trained as in DDPG
(\cref{ch:continuous}); since the critic sees all the actions, its
environment is stationary. MAPPO does the same with PPO and a centralised
value function, and is a strong baseline for cooperative tasks.

\paragraph{Value decomposition.} For a team with a common reward, a joint
Q-function $Q_{\mathrm{tot}}(s, \mathbf{a})$ cannot be maximised by
enumeration (the number of joint actions is exponential in $n$), and the
agents must act on their own. VDN (Sunehag et al., 2018) writes
$Q_{\mathrm{tot}} = \sum_i Q_i(o_i, a_i)$; QMIX (Rashid et al., 2018)
combines the $Q_i$ with a mixing network whose weights, produced by a
hypernetwork from the state, are non-negative.

\begin{proposition}{Monotonic mixing}{qmix}
  If $Q_{\mathrm{tot}}(s, \mathbf{a}) = f_s(Q_1(o_1, a_1), \dots, Q_n(o_n,
  a_n))$ with $f_s$ non-decreasing in each argument, then the joint action
  made of the individual greedy actions $a_i^\star = \argmax_{a_i} Q_i(o_i,
  a_i)$ maximises $Q_{\mathrm{tot}}(s, \cdot)$.
\end{proposition}

\begin{proof}
  For any joint action $\mathbf{a}$, $Q_i(o_i, a_i) \leq Q_i(o_i,
  a_i^\star)$ for every $i$, and replacing the arguments of $f_s$ one at a
  time by larger ones does not decrease it: $Q_{\mathrm{tot}}(s, \mathbf{a})
  \leq Q_{\mathrm{tot}}(s, \mathbf{a}^\star)$.
\end{proof}

The joint Q-function is trained by the DQN loss on the common reward, its
maximum in the target is computed agent by agent, and each agent acts
greedily on its own $Q_i$. The price is representational: monotonic
mixtures cannot express every joint Q-function (for instance, when the
best action of an agent depends on what another does).

\begin{pacman}{Learning ghosts}{ghosts-team}
  The ghosts of \code{ghost\_agents.py} follow fixed rules. Making them
  learn turns Pacman into a two-team zero-sum Markov game: Pacman against
  a cooperative team of ghosts with the opposite reward. Training them with
  QMIX or MAPPO against a frozen PPO Pacman, then retraining Pacman against
  the new ghosts, is self-play between teams, with the risk of cycling
  discussed above; a league of past ghosts and past Pacmans is the standard
  remedy.
\end{pacman}

\section{Summary}

\begin{itemize}[nosep]
  \item Markov games generalise MDPs to several agents; the solution
    concept is an equilibrium, not an optimal policy.
  \item In two-player zero-sum games, Shapley's operator is a contraction
    and minimax-Q learns the value; self-play, fictitious play and leagues
    train agents against themselves.
  \item For cooperative teams, centralised training with decentralised
    execution (MADDPG, MAPPO) restores stationarity, and monotonic value
    decomposition (VDN, QMIX) makes the joint maximisation decentralised.
\end{itemize}
